\begin{statementbox}
	\textbf{\textcolor{CStatement}{条件}}
	\begin{description}[style=nextline, leftmargin=2.8em, labelsep=0.5em, itemsep=0.35em]
		\item[\textbf{\textcolor{CStatement}{条件1（样本要求）：}}] 样本数据为 $(x_1,y_1),\dots,(x_n,y_n)$，其中 $n\ge 3$；散点图呈现大致线性趋势，且无明显异常值。
	\end{description}

	\textbf{\textcolor{CStatement}{结论}}
	\begin{description}[style=nextline, leftmargin=2.8em, labelsep=0.5em, itemsep=0.45em]
		\item[\textbf{\textcolor{CStatement}{结论一（相关系数定义）：}}]
			\textbf{\textcolor{CStatement}{直观描述：}}
			衡量两变量线性相关程度的指标，取值范围为 $[-1,1]$。
			\[
				r=\frac{\sum_{i=1}^{n}(x_i-\bar{x})(y_i-\bar{y})}
				{\sqrt{\sum_{i=1}^{n}(x_i-\bar{x})^2\sum_{i=1}^{n}(y_i-\bar{y})^2}},\quad r\in[-1,1].
			\]

		\item[\textbf{\textcolor{CStatement}{常见推论（完全线性相关）：}}]
			\textbf{\textcolor{CStatement}{直观描述：}}
			当所有样本点落在同一条直线上时，相关系数绝对值为 $1$。
			\[
				|r|=1 \iff y_i = a x_i + b,\; i=1,\dots,n.
			\]
	\end{description}
\end{statementbox}
